% SEGMENT OLP-0696-S01
% Part: proof-theory
% Chapter: propositions-as-types
% Section: type-preservation

\documentclass[../../../include/open-logic-section]{subfiles}

\begin{document}

\olfileid{int}{pty}{tpr}

\olsection{வகை மாறாமை}

இங்கு கருதப்பட்ட குறைப்பு உருமாற்றங்களையும் வரிசைமாற்று
உருமாற்றங்களையும் \Log{tN2i}, \Log{tN3i} முறைமைகளுக்கே
நேரடியாக வரையறுக்கலாம். \Log{tN3i} முறைமை, சூழல்~$\Gamma$
சார்ந்து $N$ என்ற $\lambd$-சொல்லின் வகையைப் பின்வருமாறு
வரையறுக்கிறது: $\Gamma \Sequent N : !A$ என்பதற்கு
\Log{tN3i} இல் !!a{proof} இருந்தால், அந்த வகை
!!{formula}~$!A$ ஆகும். $\lambd$-சொற்களுக்கான
$\beta$-குறைத்தல் விதிகள், !!{proof}s இன் குறைப்பு மற்றும்
வரிசைமாற்று உருமாற்றங்களைத் துல்லியமாகப் பிரதிபலிப்பதால்
ஒரு முக்கிய விளைவு கிடைக்கிறது: சூழல்~$\Gamma$
சார்ந்த $\lambd$-சொல்லின் வகை, $\beta$-மாற்றத்திற்குப்
பின்னும் மாறாது. இப்பண்பு \emph{வகை மாறாமை} எனப்படுகிறது.

\begin{prop}
$\Gamma$ சார்ந்து $N$ இன் வகை~$!A$ ஆகவும்
$N \redone N'$ ஆகவும் இருந்தால், $\Gamma$ சார்ந்து
$N'$ இன் வகையும்~$!A$ ஆகும்.
\end{prop}

% SEGMENT OLP-0696-S02
\begin{proof}
  $N$ மீதும், $\Gamma \Sequent N : !A$ என்பதன்
  !!{proof}~$\delta$ இன் உயரம் மீதும் தொகுத்தறிந்து
  பின்வருவதை எளிதில் நிறுவலாம். $M$ என்பது~$N$ இன்
  ஓர் உட்சொல்லாக இருந்தால், $\Gamma' \Sequent M : !B$
  இல் முடியும் துணை-!!{proof}~$\delta_1$ ஒன்று
  $\delta$ வில் இருக்கும். $M$ ஒரு ரெடெக்ஸ் என்றால்,
  $\delta_1$ க்கு ஓர் உருமாற்றம் பொருந்தும்.
  $\delta_1$ க்குக் குறைப்பைப் பயன்படுத்தி,
  $\Gamma' \Sequent M' : !B$ என்பதன்
  !!a{proof}~$\delta_1'$ ஐப் பெறுகிறோம்.
  \Log{tN3i} இன் உருமாற்றங்களைப் பரிசோதித்தால்,
  $M'$ என்பது~$M$ இன் சுருக்கவிளைவு என்று தெரிகிறது.
  $\delta$ வில் $\delta_1$ ஐ $\delta_1'$ ஆல்
  மாற்றுவதால், $\Gamma \Sequent N' : !A$ என்பதன்
  !!a{proof}~$\delta'$ கிடைக்கிறது; இங்கு $N$ இல்
  உட்சொல்~$M$ ஐ~$M'$ ஆல் மாற்றியதே~$N'$.

% SEGMENT OLP-0696-S03
  எடுத்துக்காட்டாக, \Intro{\land} ஐத் தொடர்ந்து
  \Elim{\land} வரும் குறைப்பு உருமாற்றத்தைப் பார்ப்போம்:
  \[
  \bottomAlignProof
  \AxiomC{}
  \Deduce$\Gamma \fCenter N_1 : !B_1$
  \AxiomC{}
  \Deduce$\Gamma \fCenter N_2 : !B_2$
  \RightLabel{\Intro{\land}}
  \BinaryInf$\Gamma \fCenter \pair{N_1}{N_2} : !B_1 \land !B_2$
  \RightLabel{\Elim{\land}[i]}
  \UnaryInf$\Gamma \fCenter \proj{i}{\pair{N_1}{N_2}} : !B_i$
  \DisplayProof
  \quad\redone\quad
  \bottomAlignProof
  \AxiomC{}
  \Deduce$\Gamma \fCenter N_i : !B_i$
  \DisplayProof
  \]
  வலப்புற !!{proof}~$\delta_1'$ இன் நிறுவல் சொல்
  $N_i$ என்பதைப் பார்க்கிறோம். இடப்புற !!{proof}~$\delta_1$
  இன் நிறுவல் சொல் $\proj{i}{\pair{N_1}{N_2}}$ மீது
  $\beta$-மாற்றத்தைப் பயன்படுத்தினால் கிடைப்பதும்
  துல்லியமாக இதுவே.

% SEGMENT OLP-0696-S04
  அல்லது \Intro{\lif} ஐத் தொடர்ந்து வரும்
  \Elim{\lif} அனுமானத்தையும் அதற்குப் பயன்படுத்தப்படும்
  குறைப்பு உருமாற்றத்தையும் பார்ப்போம்:
  \begin{multline*}
  \bottomAlignProof
  \AxiomC{$\delta_2$}
  \Deduce$x: !B_1, \Gamma \fCenter N : !B_2$
  \RightLabel{\Intro{\lif}}
  \UnaryInf$\Gamma \fCenter \lambd[\typeof{x}{!B_1}][N] : !B_1 \lif !B_2$
  \AxiomC{}
  \RightLabel{$\delta_3$}
  \Deduce$\Gamma \fCenter M : !B_1$
  \RightLabel{\Elim{\lif}}
  \BinaryInf$\Gamma \fCenter (\lambd[\typeof{x}{!B_1}][N] M) : !B_2$
  \DisplayProof
  \quad\redone\\
  \bottomAlignProof
  \AxiomC{}
  \RightLabel{$\delta_3$}
  \Deduce$\Gamma \fCenter M : !B_1$
  \RightLabel{$\Subst{\delta_2}{\delta_3}{x:!B_1}$}
  \Deduce$\Gamma \fCenter \Subst{N}{M}{x} : !B_2$
  \DisplayProof
  \end{multline*}
  இங்கும் வலப்புற !!{proof} இன் நிறுவல் சொல்,
  இடப்புற !!{proof} இன் நிறுவல் சொல்லை
  $\beta$-குறைத்தால் கிடைக்கும் சொல்லே.
  நிச்சயமாக, $\delta_2$ இன் உயரம் மீது
  தொகுத்தறிந்து ஒரு விவரத்தைச் சரிபார்க்க வேண்டும்:
  $\delta_2$ வில் உள்ள
  $\Gamma' \Sequent x:!B_1$ வடிவிலான எல்லா
  அடிக்கோள்களையும், $\Gamma \Sequent M:!B_1$
  என்பதன் !!{proof}~$\delta_3$ ஆல் மாற்றினால்,
  உண்மையில் $\Gamma \Sequent
  \Subst{N}{M}{x} : !B_2$ என்பதன்
  !!a{proof}~$\Subst{\delta_2}{\delta_3}{x:!B_1}$
  கிடைக்கிறது என்பதைச் சரிபார்க்க வேண்டும்.
\end{proof}

% SEGMENT OLP-0696-S05
!!{proof}s க்கும் $\lambd$-சொற்களுக்கும் இடையேயான
நெருங்கிய பொருத்தத்திற்கு இன்னொரு விளைவு உண்டு.
சூழல்~$\Gamma$ சார்ந்து வகை~$!A$ உடைய
$\lambd$-சொல்~$N$ இல் ரெடெக்ஸ்~$M$ இருந்தால்
(அதாவது, அது இயல்புவடிவத்தில் இல்லாவிட்டால்),
$\Gamma \Sequent N : !A$ என்பதற்குரிய
\Log{tN3i} நிறுவல்~$\delta$ வில்,
$\Gamma' \Sequent M: !B$ இல் முடியும்
உருமாற்றம் பொருந்தக்கூடிய ஒரு துணை-!!{proof}
இருக்கும். $M$ ஐ அதன் சுருக்கவிளைவால் மாற்றி
$N \redone N'$ என்று பெற்றால்,
$\delta$ வில் அதற்குரிய குறைப்பைப் பயன்படுத்துவதன்
விளைவு $\Gamma \Sequent N' : !A$ என்பதன்
!!a{proof} ஆகும். ஆகவே, இயல்புவடிவத்தில் இல்லாத
ஒவ்வொரு $\lambd$-சொல்லும் $\beta$-குறைத்தல்
மூலம் மற்றொரு சொல்லுக்குச் செல்லும்.
இப்பண்பு \emph{முன்னேற்றம்} எனப்படுகிறது.

முன்னேற்றமும் வகை மாறாமையும் சேர்ந்து,
$\lambd$-சொற்களின் $\beta$-குறைத்தல்
“தேங்கி நிற்காது” என்பதை உறுதிசெய்கின்றன:
ஒரு சொல் சரியாக வகையிடப்பட்டு,
இயல்புவடிவத்தில் இல்லாவிட்டால்,
அது அதே வகையுடைய மற்றொரு சொல்லாகக் குறையும்.

\end{document}
